\section{Test environments}
\label{sec:env}

All experiments use spatially cued orientation-change-detection tasks built to
reproduce the essential structure of the paradigms used to study covert attention
in non-human primates \citep{posner1980attention, cohen2009attention,
herman2017color, srinath2021attention}. The agent views a sequence of
$50\times50$ frames and, on each frame, chooses to \emph{wait} or to \emph{declare
a change}. The design isolates internal attentional state from sensory processing
by separating the cue in time from the change it predicts \citep{posner1980attention,
muller1987spatial}, and it makes attentional engagement parametric through cue
\emph{reliability} (validity) \citep{mayo2016graded}.

\subsection{Base task: cued change detection with value}
\label{sec:env-base}
Each trial unfolds over a fixed number of frames. The trial opens with a blank
frame; a spatial cue is then presented briefly and removed, after which the display
is blank again before the stimulus array appears; the array of oriented
``Gabor''-like patches is then shown for the remainder of the trial. On half of
trials, at a change time that is \emph{jittered across frames}, the orientation of
one patch steps by a magnitude $\Delta$ and persists; on the other half no change
occurs. Frame-to-frame orientation noise is added throughout to titrate difficulty
and to require genuine evidence integration rather than single-frame comparison.

The cue carries two pieces of information. Its \emph{location} (here, a side ---
left or right of the array) indicates where a change is most likely to occur, and
its \emph{reliability} (validity) --- the probability that a change, if one occurs,
appears at the cued location --- is drawn from $\{0.25, 0.50, 0.75, 1.00\}$ and is
displayed by the completeness of a ring around the cue. The cue's \emph{colour}
sets the reward magnitude available on that trial (a value manipulation, analogous
to reward-biased attention paradigms \citep{stanisor2013neural, baruni2015reward,
maunsell2004neuronal}). Reward is delivered for a correct decision: declaring at or
after the change is a hit (rewarded by the cue-colour value); declaring before the
change is a false alarm; waiting through a no-change trial is a correct rejection
(rewarded); failing to declare a change is a miss. Because change onset is jittered
and the response is event-locked, the agent cannot exploit a fixed temporal
strategy and must instead maintain the cue and monitor for the change.

This base task is the setting for the central architectural comparison
(Section~\ref{sec:results}). It combines a \emph{validity} manipulation (graded cue
reliability) with a \emph{value} manipulation (cue colour), so that the trained
agent can be probed both for the classic reliability-scaled cueing benefit and for
value-directed responding.

\subsection{Set-size variant: nine stimuli}
\label{sec:env-setsize}
To ask how attentional benefit scales with the number of competing locations --- a
question with a long history in the human and animal literatures
\citep{palmer1995psychophysics, mazyar2012does, eckstein2013visual} --- we
constructed a set-size variant with nine stimuli arranged in a $3\times3$ grid. A
location cue marks one of the nine cells, and its ring completeness again signals
reliability, now drawn from five levels $\{0, 0.25, 0.50, 0.75, 1.00\}$, where a
proportion of $0$ is an uninformative cue (uniform over all nine locations) and a
proportion of $1$ is a perfectly valid cue. Reward is uniform across locations
(removing the value manipulation) so that any reliability-graded effect reflects
validity alone. With nine locations rather than four, knowing where to look should
be worth more: the set-size variant tests whether the cue-reliability benefit, and
the proportion-graded recruitment of attention, \emph{grow} with set size --- the
set-size $\times$ validity interaction analysed in Section~\ref{sec:setsize}.

\subsection{Distractor variant}
\label{sec:env-distractor}
The base task contains no competing events: when a change occurs it is always the
task-relevant target. Natural attention, however, is defined by the need to select
a target while \emph{suppressing} salient but irrelevant events
\citep{desimone1995neural, sprague2018rescuing, cosman2018prefrontal}. We therefore
constructed a distractor variant in which, independently of the target, the uncued
side may undergo its own orientation change with probability $p_{\mathrm{dist}}$
(with its own location, magnitude and timing). A change on the uncued side is a
distractor that must be ignored: declaring in response to it is a false alarm that
ends the trial without reward. Reward remains keyed to the target alone. The
distractor variant thus turns the task into a selective-attention problem and lets
us ask, of the split architecture, \emph{where} in the two pathways an irrelevant
event is filtered, and whether value modulates that filter
(Section~\ref{sec:distractor}).
